\section{Introduction}
The primary objective of this practical report is to investigate the relationship between Glycosolated Hemoglobin (GHB) and a set of physiological risk factors in a sample of African American subjects from central Virginia, as originally presented by Faraway \cite{farway2016}. GHB is an important marker for diabetes diagnosis, where levels above 7.0 typically indicate a positive diagnosis. Given the continuous and strictly positive nature of the response variable (GHB), combined with the right-skewness of the biological data, standard linear regression models (such as Ordinary Least Squares), which assume normality (Gaussian distribution), may be inappropriate. Therefore, this report proposes a Generalized Linear Model (GLM) assuming a Gamma distribution with a log-link function.

Our analysis follows a Bayesian perspective, utilizing Markov Chain Monte Carlo (MCMC) methods for parameter estimation, followed by model selection based on the Deviance Information Criterion (DIC). Finally, to ensure the robustness of the findings and satisfy the requirement for innovation, the results are validated using a non-parametric Bootstrap approach.